\section{模型验证与灵敏度分析}

\subsection{模型验证}

模型的可信度来自四类验证（\tabref{tab:allverify}）：解析解验证、统计性质验证、
物理机制验证与数据交叉验证。所有验证都可在仓库中一键重跑。

\begin{table}[H]
\centering\small
\caption{模型验证汇总}
\label{tab:allverify}
\begin{tabular}{@{}llll@{}}
\toprule
类别 & 验证内容 & 判据 & 结果\\
\midrule
解析 & 终端速度：解析解 vs 数值积分 & 相对误差 & 0.00\%\\
统计 & OU 过程方差 vs 理论 $\sigma^2$ & 相对误差 & 1.61\%\\
统计 & 30\,000 粒子收敛（60 案例） & $P_{72}$ 标准误 & 0.3--0.6 百分点\\
物理 & 滑移判据：单坡解析 vs 逐 20 m 重算 & 停止位置一致 & 通过\\
物理 & 海水状态方程 EOS-80 vs 公开值 & 相对误差 $<10^{-4}$ & 通过\\
数据 & ETOPO 站点深度 vs 海图 & 偏差 $<2$ m & 通过\\
数据 & HYCOM 流速量级 vs 文献 & 同量级 & 通过\\
\bottomrule
\end{tabular}
\end{table}

\subsection{灵敏度分析汇总}

\tabref{tab:sensitivity} 汇总了所有做过的灵敏度分析。按对结论的影响从大到小排序：
\textbf{响应延迟}（每 6 h 使 $P_{72}$ 下降 8--22 个百分点）$>$ \textbf{增援到位时间}
（最简方案每 3 h 下降 0.6--0.9 个百分点）$>$ \textbf{协调参数 $\beta$}
（0 到 5 使 $P_{72}$ 从 97.2\% 降到 91.8\%）$>$ \textbf{海域}
（不超过 2.4 个百分点）$>$ \textbf{蒙特卡洛误差}（0.3--0.6 个百分点）。

\begin{table}[H]
\centering\small
\caption{灵敏度分析汇总（主指标为 $P_{72}$）}
\label{tab:sensitivity}
\begin{tabular}{@{}llll@{}}
\toprule
因素 & 变化范围 & 指标变化 & 结论\\
\midrule
搜救启动延迟 & 0/6/12/24 h & M7: 95.5\% $\to$ 47.4\% & 影响最大\\
增援到位时间 & 0/3/6/12 h & 2 船: 97.5\% $\to$ 94.8\% & 次之\\
重复扫描惩罚 $\beta$ & 0--5 & 97.2\% $\to$ 91.8\% & 可调，取 1.0\\
海域 & 地中海/加勒比 & $\pm$2.4 个百分点 & 结论稳健\\
扫宽 $W$ & $\pm$30\%（解析） & 方案排序不变 & 结论稳健\\
装备速度 & 给定表（kn） & 已含 & —\\
蒙特卡洛误差 & 60 案例 & $\pm0.3$--$0.6$ & 可接受\\
\bottomrule
\end{tabular}
\end{table}

\subsection{局限性与未来工作}

\begin{enumerate}[leftmargin=1.8em,itemsep=0pt]
\item \textbf{完美预报假设}：预测与真值共用同一套 HYCOM 流场，未建模预报误差。
      这是最重要的局限；改进方向是引入集合预报（如 HYCOM 集合或加扰动的随机流场），
      让搜救方只掌握“预报 + 误差”；
\item \textbf{探测模型参数}：扫宽 $W$ 是量级估计，未经海试标定；水柱探测的
      POD 与目标姿态、深度、海况强相关，未来应做装备性能试验；
\item \textbf{失效概率为专家判断}：S0--S4 的概率、浮力异常分布都无公开统计，
      只能用敏感性分析覆盖；
\item \textbf{未建模声信标/ULB}：若潜水器携带 37.5 kHz USBL/声信标，被动监听可把
      搜救问题从“扫海”变成“定位”，效率高一个量级——这是最有价值的改进方向之一；
\item \textbf{未建模多目标}：多艘潜水器同时失事时的分配与协调未涉及；
\item \textbf{未建模装备可靠性/海况限制}：AUV 续航、回收周期、船间通信延迟、
      恶劣海况下的降级都未纳入。
\end{enumerate}
